%% THERMODYNAMIQUE %%
%
\section{Rapport de m\'elange en vapeur d'eau \`a partir de l'humidit\'e absolue (approximation) }

\itemstyle{ Op\'erateur~: } {\sf RapM\_Ha\_approx}
\index{RapM\_Ha\_approx}

%%%%%OBJET%%%%%
\itemstyle{ Objet~: }
Calcul approch� du rapport de m\'elange en vapeur d'eau (rapport entre la
masse de vapeur d'eau et la masse d'air \textbf{sec}, en g/kg)
\`a partir de l'humidit\'e absolue de l'air (en g/m$^3$).

Ce calcul est une division par la masse volumique ("densit�") de l'air (kg/m3).

L'approximation r�side dans le fait qu'on fait ce calcul avec
les constantes de l'air sec (Ra). Cette approximation peut �tre
justifi�e dans les cas o� l'humidit� est faible (haute atmosph�re)
ou que l'on part d'une temp�rature sonique, voir virtuelle (dans ce cas,
le calcul est exact).

%%%%%ARGUMENTS%%%%%
\bigskip
%%%%%ENTREES%%%%
\begin{tabarg}
\multicolumn{3}{l}{\titretab{Entr\'ees}}\\
\hline
$H_a$ & Vecteur, {\tt float} & Humidit\'e absolue ($g/m^3$)\\
$P_s$ & Vecteur, {\tt float} & Pression statique ($hPa$)\\
$T_s$ & Vecteur, {\tt float} & Temp\'erature statique ($K$ ou $^\circ C$)\\
\end{tabarg}

\medskip
%%%%%SORTIES%%%%
\begin{tabarg}
\multicolumn{3}{l}{\titretab{Sorties}}\\
\hline
$RapM$ & Vecteur, {\tt float} & Rapport de m\'elange en vapeur d'eau ($g/kg$)\\
\end{tabarg}


%%%%%FORMULE%%%%
\itemstyle{ Formule }

%%% THERMODYNAMIQUE
\[
RapM = \frac{ R_a \times T_s }{P_s \times 100.} . H_a
\]
avec $R_a = 287.05 J.kg^{-1}.K^{-1} $, et $T_s$ exprim\'e en K.


%%%%%AUTEUR%%%%%
\itemstyle{ Auteur~:}
CNRM/GMEI/TRAMM

%%%%%REFERENCES%%%%%
\itemstyle{ R\'ef\'erences~:}
Triplet-Roche\cite{Triplet-Roche}\\
D�finitions des humidit�s (relatives, absolues, sp�cifiques)
selon les tables m�t�orologiques internationales
de l'OMM \cite{OMM-188.TP.94}.
